\documentclass{report}
\usepackage[main=brazilian, english]{babel}
\usepackage{hyperref}

\title{Implementação de um bate-papo utilizando o protocolo MQTT na linguagem Java}
\author{Erickson G. Müller}
\date{11 de Novembro de 2026}

\begin{document}
\maketitle

\chapter{Introdução}

O projeto de bate-papo com o protocolo MQTT foi desenvolvido em Java utilizando a biblioteca Paho para essa linguagem \footnote{\url{https://eclipse.dev/paho/files/javadoc/index.html}} e o broker Mosquitto \footnote{\url{https://mosquitto.org/}}. A interface é manuseada por terminal através de comando indicados pelo próprio programa. O trabalho foi desenvolvido individualmente com o auxílio do modelo de linguagem Gemini Flash Lite de forma não-agêntica.

A proposta do trabalho foi desenvolver um sistema de bate-papo um-a-um e em grupo, possibilitando a listagem de usuários disponíveis e ausentes, listagem de grupos e seus respectivos líderes e membros, gerenciamento de grupos, notificações e histórico de ações no sistema de maneira cronológica.

\chapter{Instalação}
Para compilar e rodar a aplicação, é necessário ter o \textbf{Java Development Kit (JDK)} e o broker \textbf{Eclipse Mosquitto} instalados. Certifique-se de que a biblioteca Paho (arquivo `.jar`) está localizada dentro da pasta `lib/` do projeto.

\section{Dependências}
\begin{verbatim}
Linux:
sudo pacman -S jdk-openjdk mosquitto
ou
sudo apt update
sudo apt install default-jdk mosquitto

MacOS:
brew install java mosquitto

Windows:
winget install Microsoft.OpenJDK.17
\end{verbatim}

\section{Compilação}
A compilação do programa deve ser executada com o seguinte comando:
\begin{verbatim}
Linux e MacOS:
javac -cp "lib/*" -d bin src/chatmqtt/*.java

Windows:
javac -cp "lib\*" -d bin src\chatmqtt\*.java
\end{verbatim}

\section{Inicialização do Broker}
Após as instalações necessárias, o programa necessita que você inicie o broker local para rotear as mensagens. Abra um terminal exclusivo e digite o seguinte comando:

\begin{verbatim}
Linux e MacOS:
mosquitto -v

Windows:
mosquitto.exe -v
\end{verbatim}

\section{Execução}
Para a execução de cada cliente, inicie o programa com o seguinte comando:
\begin{verbatim}
Linux e MacOs:
java -cp "bin:lib/*" chatmqtt.chat

Windows: 
java -cp "bin;lib\*" chatmqtt.chat
\end{verbatim}

\chapter{Arquitetura e Classes}
O software foi desenvolvido com uma arquitetura modularizada, separando as partes de comunicação de rede, interface e lógica de estados. A aplicação é dividida em seis classes principais:

\section{\texttt{controller.java}}
Esta classe atua como a ponte de comunicação entre a aplicação e o broker Mosquitto, encapsulando os métodos da biblioteca Eclipse Paho MQTT. Suas principais responsabilidades são:
\begin{itemize}
    \item Estabelecer a conexão TCP persistente com o servidor (utilizando a flag \texttt{setCleanSession(false)}).
    \item Assinar e desassinar tópicos conforme a criação e encerramento de conversas, para isso foram criados os métodos \texttt{assinarTopico} e \texttt{desassinarTopico}.
    \item Publicar mensagens padrão (com QoS 1) e mensagens retidas (para controle de status e registro em grupos).
    \item Implementar o \textit{callback} assíncrono com o método principal (\texttt{messageArrived}), que atua como um roteador central. Ele intercepta todos os pacotes recebidos pela rede e os despacha para as classes de gerenciamento adequadas com base no prefixo do tópico.
\end{itemize}

Alguns métodos foram criados para automatizar a configurações que se repetem no gerenciamento do callback, como \texttt{desconectar(), enviarMensagem(), enviarMensagemRetida() e removerMensagemRetida()}.

\begin{verbatim}
	public void publicarStatus(String status) throws MqttException {
		String messageStatus = userId + ":" + status;
		MqttMessage statusUser = new MqttMessage(messageStatus.getBytes());
		statusUser.setQos(1);
		statusUser.setRetained(true);// para quem logar depois
		client.publish("USERS/" + userId, statusUser); // Publica no subtópico
	}
	...
				public void messageArrived(String topic, MqttMessage message) throws Exception {
				String payload = new String(message.getPayload());

				if (topic.startsWith("USERS/")) {// Criar subtópicos para evitar o problema de
									// Guardar apenas uma mensagem
					gerenciadorUsuarios.atualizarStatus(payload);
				}
				else if (topic.startsWith("GROUPS/")){ // Cria subtópico GROUPS/
					gerenciadorGrupos.atualizarGrupo(topic, payload);
				}
	...

\end{verbatim}

\section{\texttt{session.java}}
É o motor central de gerenciamento de estado das conversas e interações do usuário. A classe \texttt{session} não interage diretamente com o MQTT, mas mantém toda a memória operacional da aplicação. Ela gerencia:
\begin{itemize}
	\item \textbf{Histórico e Conversas Ativas:} Mantém na memória todas as salas de bate-papo abertas e as mensagens trocadas. \texttt{private Map<String, List<String>> historicoMensagens;} e \texttt{private Map<String, String> conversasAtivas;}
    \item \textbf{Controle de Tela :} Armazena a referência da ``sala ativa'' atual na tela do terminal. Se uma mensagem chega para a sala que o usuário está visualizando, ela é impressa instantaneamente; caso contrário, a classe emite uma notificação assíncrona não obstrutiva.
    \item \textbf{Solicitações:} Controla as filas (utilizando estruturas \texttt{HashSet} e \texttt{HashMap}) e os logs de pedidos pendentes de conversas privadas e de ingressos em grupos.
\end{itemize}
\begin{verbatim}
	public void registrarMensagemChat(String topico, String mensagem, String userId) {
	...
		if (topico.equals(salaAtivaNaTela)) {
			if (!mensagem.startsWith("["+userId+"]")) {
				System.out.println(mensagem);
			}
		}
		else {
			if (!mensagem.startsWith("["+userId+"]")) {
				System.out.print("\r\n[NOTIFICAÇÃO] Nova mensagem recebida\n>> ");
			}
		}
...
	public void listarPendentes() {
		System.out.println("\n--- Solicitações de Conversa ---");
		if (solicitacoesPendentes.isEmpty()) {
			System.out.println("Nenhuma solicitação pendente");
		}
		else{
			for (String solicitante : solicitacoesPendentes){
				System.out.println("- " + solicitante);
			}
		}
	}
	public void listarLogSolicitacoes() {
		System.out.println("Registros de solicitações em ordem cronológica");
		if (logSolicitacoes.isEmpty()){
			System.out.println("Nenhuma solicitação registrada");
		}
		else{
			for (String registro : logSolicitacoes){
				System.out.println(registro);
			}
		}
	}
	public void registrarLogSolicitacoes(String tipo, String remetente, String destinatario,
	String canal){
		String timestamp = LocalDateTime.now().format(formatter);
		String log = "["+tipo+"]"+timestamp+" " + remetente + "->" + destinatario;
		if (tipo.equals("ACEITO")){
			log += " ("+canal+")";
		}
		logSolicitacoes.add(log);
	}
	public void solicitarParticipacao(String membro, String grupo){
		solicitacoesGrupo.put(membro,grupo);
		System.out.println("\r\n[NOTIFICAÇÃO] Nova solicitação de participação ao grupo
		'"+grupo+"' vindo de "+membro );
		System.out.print(">> ");
	}

	public void removerSolicitacaoGrupo(String membro){
		solicitacoesGrupo.remove(membro);
	}
	public boolean possuiSolicitacaoGrupo(String membro){
		return solicitacoesGrupo.containsKey(membro);
	}
	public String getGrupoSolicitado(String membro){
		return solicitacoesGrupo.get(membro);
	}
	public void listarSolicitacoesPendentesGrupo(){
		System.out.println("Solicitações de participação Pendentes:");
		if (solicitacoesGrupo.isEmpty()){
			System.out.println("Nenhuma solicitação no momento");
			return;
		}

		for (Map.Entry<String, String> entry : solicitacoesGrupo.entrySet()){
			System.out.println("Usuário '"+entry.getKey()+"' solicitou participação no grupo '"
			+entry.getValue()+"'");
		}
	}
\end{verbatim}

\section{\texttt{groups.java}}
Classe voltada exclusivamente à manutenção das estruturas de dados das salas de chat coletivas. O sistema utiliza esta classe para recriar a estrutura hierárquica dos grupos na memória RAM de cada cliente. A classe mantém um dicionário (\texttt{Map}) que vincula o nome de cada grupo a um objeto interno chamado \texttt{infoGroup}. Este objeto registra estritamente qual é o ID do líder e um conjunto de dados (\texttt{Set}) com os IDs de todos os membros comuns, permitindo a validação rápida de permissões (como quem pode deletar o grupo ou enviar mensagens).

\section{\texttt{users.java}}
Responsável pelo controle de presença da aplicação. Processa as mensagens retidas vindas do tópico raiz \texttt{USERS/\#} e atualiza instantaneamente a lista local informando quais clientes estão \texttt{ONLINE} ou \texttt{OFFLINE}. Isso garante que o menu de listagem de usuários reflita o estado real e sincronizado da rede em qualquer instante.

\section{\texttt{menu.java}}
Implementa a Interface de Linha de Comando (CLI) que guia a experiência visual e a interatividade do usuário com o sistema. Suas funções abrangem:
\begin{itemize}
    \item Capturar as entradas de teclado (através da classe \texttt{Scanner}) e processar a navegação nos submenus hierárquicos de conversas e grupos.
    \item Bloquear o fluxo contínuo de execução durante um bate-papo ativo (através de laços de repetição \texttt{while}) para permitir o envio de mensagens de texto sem a interferência do menu principal.
    \item Intermediar a lógica complexa de saída, abandono e exclusão de conversas e grupos, chamando coordenadamente os métodos de limpeza visual da classe \texttt{session} e as funções de exclusão de mensagens retidas da classe \texttt{controller}.
\end{itemize}

\chapter{Funcionalidades}
\section{Comunicação um-a-um}
Quando um usuário deseja iniciar um bate-papo, ele envia uma solicitação que fica pendente até o aceite do destinatário. Após a aceitação, o sistema gera dinamicamente um tópico exclusivo para a sessão (ex: \texttt{Usuario1\_Usuario2\_Timestamp}), garantindo o isolamento da conversa. As mensagens de texto são enviadas com Qualidade de Serviço 1 (QoS 1 - \textit{At least once}) e a conexão é configurada com \texttt{setCleanSession(false)}, o que permite que o broker Mosquitto enfileire mensagens caso um dos usuários perca a conexão ou fique offline temporariamente, entregando-as assim que a rede for restabelecida. O chat possui auto-recuperação de estado e um comando \texttt{/encerrar} para destruir a sessão e desassinar o tópico mutuamente.
\subsection{Listagem de Usuários e Status}
Para o controle de presença (online/offline), a aplicação utiliza mensagens retidas (\textit{Retained Messages}) do protocolo MQTT. Assim que um cliente conecta ao broker, ele publica seu status no tópico \texttt{USERS/SeuID} com a flag de retenção ativada e carga útil (\textit{payload}) de \texttt{ONLINE}. Na desconexão, o status \texttt{OFFLINE} é emitido. Como as mensagens são retidas, qualquer usuário que entrar no sistema posteriormente recebe instantaneamente o panorama atualizado de quem está disponível, gerenciado localmente via estruturas de dados \texttt{HashSet}.

\section{Comunicação em grupo}
O sistema permite a criação de salas de bate-papo coletivas de forma descentralizada. O usuário que cria o grupo é designado automaticamente como Líder. Usuários comuns podem listar os grupos existentes e enviar solicitações de entrada. O líder possui autonomia para aprovar ou recusar pedidos e é o único com permissão para dissolver o grupo (eliminando a sala para todos os participantes). Membros comuns podem abandonar o grupo a qualquer momento através do comando \texttt{/abandonar} ou através do sub-menu de grupos. A troca de mensagens ocorre em um tópico dedicado com o prefixo \texttt{CHAT\_GRUPO/}.
\subsection{Listagem de Grupos}
Assim como na listagem de usuários, a estruturação dos grupos tira proveito das mensagens retidas. A participação é registrada no broker sob o padrão \texttt{GROUPS/NomeDoGrupo/Usuario}, com a carga útil definindo o nível de permissão (\texttt{LEADER} ou \texttt{MEMBER}). Para remover um membro ou apagar um grupo inteiro, o sistema envia uma mensagem com \textit{payload} vazio (zero bytes) para o respectivo tópico, sinalizando ao Mosquitto que a mensagem retida deve ser deletada, o que atualiza as listas locais de todos os clientes conectados.

\section{Organização de Tópicos}
Para manter a segregação lógica e evitar o cruzamento de mensagens de dados com mensagens de controle, a árvore de tópicos foi estruturada da seguinte maneira:
\subsection{Tópico de controle do cliente}
Cada usuário assina obrigatoriamente um canal particular com o sufixo \texttt{\_Control} (ex: \texttt{Erickson\_Control}). Este tópico é oculto e dedicado exclusivamente à sinalização do sistema. É através dele que transitam requisições (\texttt{REQ}, \texttt{GROUP\_REQ}), aprovações (\texttt{ACCEPT}, \texttt{GROUP\_ACCEPT}) e rejeições, funcionando como uma camada de controle separada da camada de texto dos chats.
\subsection{Tópico de controle de grupos}
O gerenciamento dos grupos é dividido em dois ramos principais:
\begin{itemize}
    \item \texttt{GROUPS/\#}: Tópico de administração assinado por todos os usuários de forma global. Usado estritamente para manter a lista de grupos, líderes e membros sincronizada na memória RAM de cada cliente.
    \item \texttt{CHAT\_GRUPO/\#}: Tópicos de comunicação efetiva. Apenas os membros validados de um grupo assinam o respectivo canal de texto, garantindo privacidade e evitando o recebimento de mensagens de grupos dos quais o usuário não faz parte.
\end{itemize}

\section{Notificações}
O recebimento de mensagens e eventos operacionais baseia-se num modelo totalmente assíncrono utilizando \textit{Callbacks} (\texttt{messageArrived}) da biblioteca Eclipse Paho. Isso permite que a aplicação rode continuamente, aguardando interações do usuário no console principal, enquanto uma \textit{thread} secundária intercepta pacotes do broker e imprime notificações de mensagens recebidas, pedidos de amizade e alertas de sistema na tela de forma não obstrutiva.
\begin{verbatim}
void setCallback(MqttCallback callback)

Sets the callback listener to use for events that happen asynchronously.

void messageArrived(java.lang.String topic,
                    MqttMessage message)
             throws java.lang.Exception
\end{verbatim}


\end{document}
